% 半人马小行星
% license CCBYSA3
% type Wiki

本文根据 CC-BY-SA 协议转载翻译自维基百科\href{https://en.wikipedia.org/wiki/Centaur_(small_Solar_System_body)}{相关文章}。

\begin{figure}[ht]
\centering
\includegraphics[width=6cm]{./figures/4207f1463529dc38.png}
\caption{} \label{fig_Centau_2}
\end{figure}
在行星天文学中，半人马天体（centaur）是指一类在木星与海王星之间绕太阳运行、并且其轨道会与一颗或多颗巨行星轨道相交的小型太阳系天体。由于这一特性，半人马天体的轨道通常是不稳定的；几乎所有半人马天体的轨道动力学寿命都只有数百万年$^{[1]}$，但有一个已知的例外，即 514107 Kaʻepaokaʻāwela，其轨道可能是稳定的（尽管是逆行轨道）$^{[2]}$（见注释$^{[1]}$）。半人马天体通常同时表现出小行星和彗星的特征。它们的名称来源于神话中的半人马，即兼具马与人形态的生物。由于观测上对大尺寸天体存在偏向性，半人马天体的总体数量难以确定。对直径大于 1 km 的半人马天体数量估计范围从约 44,000 个$^{[1]}$到超过 1,000 万个$^{[4][5]}$不等。

按照喷气推进实验室（JPL）的定义（也是本文采用的定义），第一个被发现的半人马天体是 944 Hidalgo，发现于 1920 年。然而，直到 1977 年发现 2060 Chiron 之后，半人马天体才被视为一个独立的天体族群。目前已确认的最大半人马天体是 10199 Chariklo，其直径约为 250 km，大小相当于一颗中等尺寸的主带小行星，并且已知其拥有一套环系统。该天体发现于 1997 年。

目前尚无任何半人马天体被近距离拍摄过，不过有证据表明，土星的卫星菲比（Phoebe）——由卡西尼号探测器于 2004 年拍摄——可能是一颗起源于柯伊伯带并被俘获的半人马天体$^{[6]}$。此外，哈勃空间望远镜也获取了一些关于 8405 Asbolus 表面特征的信息。

谷神星可能起源于外行星区域$^{[7]}$，如果这一假设成立，它可以被视为一颗前半人马天体，但目前观测到的半人马天体均起源于其他区域。

在已知占据半人马式轨道的天体中，约有 30 个显示出类似彗星的尘埃彗发，其中 2060 Chiron、60558 Echeclus 和 29P/Schwassmann–Wachmann 1 在完全位于木星轨道之外的情况下仍表现出可探测的挥发物释放$^{[8]}$。因此，Chiron 和 Echeclus 同时被归类为半人马天体和彗星，而 Schwassmann–Wachmann 1 一直被归类为彗星。其他半人马天体（如 52872 Okyrhoe）也被怀疑曾出现过彗发。理论上，任何被扰动并足够接近太阳的半人马天体，最终都可能演化为彗星。
\subsection{分类}
若一颗天体的近日点距离或半长轴位于外行星之间（即木星与海王星之间），即可被归类为半人马天体。由于该区域内轨道在长期尺度上的固有不稳定性，即便是当前并未与任何行星轨道相交的半人马天体（如 2000 GM137 和 2001 XZ255），其轨道也会逐渐演化，并最终受到扰动，从而开始与一颗或多颗巨行星的轨道相交$^{[1]}$。一些天文学家仅将半长轴位于外行星区域的天体归为半人马天体；而另一些学者则认为，只要近日点位于该区域内即可归类为半人马天体，因为这两类轨道在动力学上同样是不稳定的。
\subsubsection{判据差异}
然而，不同机构在对边界天体进行分类时采用了不同的判据，这些判据通常基于轨道要素的具体取值：
\begin{itemize}
\item 小行星中心（Minor Planet Center，MPC）将半人马天体定义为近日点位于木星轨道之外（5.2 AU < q），且半长轴小于海王星半长轴（a < 30.1 AU）的天体$^{[9]}$。MPC 有时会将半人马天体与散射盘天体列为同一类群$^{[10]}$。
\item 喷气推进实验室（Jet Propulsion Laboratory，JPL）采用了相近的定义，将半人马天体界定为半长轴 a 介于木星与海王星之间的天体（5.5 AU ≤ a ≤ 30.1 AU）$^{[11]}$。
\item 相比之下，深黄道巡天（Deep Ecliptic Survey，DES）使用的是一种动力学分类方案。该分类基于将当前轨道向未来延拓 1000 万年时，其轨道行为在数值模拟中的变化。DES 将半人马天体定义为非共振天体，其瞬时（密切）近日点在模拟过程中任意时刻小于海王星的瞬时半长轴。该定义旨在与“行星交叉轨道”等价，并暗示这类天体在当前轨道上的寿命相对较短$^{[12]}$。
\item 《海王星之外的太阳系》（The Solar System Beyond Neptune，2008）一书将半长轴位于木星与海王星之间、且相对于木星的 Tisserand 参数大于 3.05 的天体定义为半人马天体；将相对于木星的 Tisserand 参数低于该阈值、并且为排除柯伊伯带天体而额外施加一个任意的近日点限制（q ≤ 7.35 AU，位于通往土星距离的一半处）的天体归类为木星族彗星；而将半长轴大于海王星、且处于不稳定轨道上的天体归类为散射盘天体$^{[13]}$。
\item 还有一些天文学家倾向于将半人马天体定义为：不与行星处于共振状态、近日点位于海王星轨道以内、并且可以证明在未来 1000 万年内很可能进入某一气体巨行星希尔球范围的天体$^{[14]}$。在这种定义下，半人马天体可被视为向内散射的天体，其与行星的相互作用更强、轨道演化也比典型的散射盘天体更快。
\item JPL 小天体数据库目前列出了 910 颗半人马天体$^{[15]}$。此外，还有 223 个跨海王星天体（即半长轴大于海王星，30.1 AU ≤ a），其近日点距离小于天王星轨道（q ≤ 19.2 AU）$^{[16]}$。
\end{itemize}
